\subsubsection{Explications TAD}
\paragraph{TAD Case}
Le TAD Case a été créé afin de pouvoir uniformiser l'accès aux données des cases voisines.

Celui-ci permet ainsi de connaitre les différentes cases qui lui sont accessibles, et avec leur direction.

Afin de ne pas surcharger les fonctionnalités, nous avons atomisé les actions possibles.
De ce fait, nous avons donc \verb|obtenirCaseHaut|, \verb|obtenirCaseBas|, \verb|obtenirCaseGauche| et \verb|obtenirCaseDroite|.
C'est ainsi une simple interface d'accès aux données stockées.
Aucune réelle logique métier supplémentaire est ajouté ici.

\paragraph{TAD Labyrinthe}

Le labyrinthe est un regroupement de cases.
Il est en quelque sorte une surcouche de logique métier au TAD Case.
Il contient cependant une information supplémentaire, celle des cases d'entrée et de sortie.

La valeur ajoutée au TAD Case qu'il englobe est la possibilité de récupérer tout les voisins d'une case.
L'algorithme de recherche de plus court chemin a été séparé du TAD Labyrinthe, puisqu'il n'a pas été vu comme indispensable à une existence d'un labyrinthe.


